\chapter{Problem Definition and Study Design}
\label{chap:problem-definition-study-design}

This chapter defines CliniCause as a framework for constructing two dataset-specific observational causal-analysis testbeds.  Each resource packages source-observed ICU records and outcomes together with an explicit analytical representation, project-specified causal assumptions, estimator interfaces, diagnostics, and provenance.  MIMIC-III and PhysioNet 2012 are instantiated separately: their shared contract identifies corresponding analytical roles, but it neither pools records nor makes similarly named proxies, graphs, adjustment sets, or estimands identical.

\section{Dataset-Specific Observational Testbed Contract}
\label{sec:problem-definition-study-design:testbed-contract}

For dataset \(d\), the constructed resource links a source-specific record identity; source-observed longitudinal measurements; observation and missingness information; a deterministic proxy-state representation; model-generated proxy annotations where used; fixed aggregation outputs where used; source-observed in-hospital mortality; observed covariates; a project-specified DAG; exposure-specific adjustment sets; an analytical cohort definition; and provenance records that distinguish the evidence class of each component.  The resource is an observational \emph{testbed} when used to examine causal-estimator behavior.  It is not a benchmark with a known counterfactual answer.

These components do not have a common evidential status.  Record identifiers, recorded measurement events, their availability patterns, observed covariates, and mortality originate in the source data subject to extraction and data-quality limitations.  Proxy definitions, cutoffs, temporal rules, cohort definitions, and covariate roles are representation-defined research choices.  Rule-derived proxy labels are deterministic derived objects instantiated from those choices and the available observations.  Learned proxy annotations are model-generated estimates of the rule-derived labels, while the archived majority-vote interface is a fixed derived analytical output.  DAGs and exposure-specific adjustment sets are project-specified causal assumptions.  Provenance and evidence-class records link these objects without implying that source observation, deterministic derivation, model estimation, and causal assumption provide equivalent evidence.

The controlled intervention in this construction therefore occurs at the representation layer.  Researchers specify which observed evidence defines a proxy, which records form an analytical cohort, which variables serve as covariates and outcomes, and which graph and adjustment assumptions govern each exposure analysis.  Deterministic source code then instantiates the selected proxy rules on patient records.  CliniCause does not simulate those patient records, a treatment or exposure-assignment mechanism, or mortality outcomes, and deterministic proxy-label construction does not make the underlying clinical state known.  The causal answer remains unknown because the data and outcomes are observational and the representation choices do not establish identification.

\section{Units, Time Horizons, and Data Objects}
\label{sec:problem-definition-study-design:units-time-horizons-data-objects}
% Section ID: C3.1.
% Evidence map: see thesis-writing/planning/chapter_evidence_map.md.

This thesis studies retrospective intensive-care records represented as irregular clinical time series.  The operational study unit is a time-series record or ICU stay, depending on the source dataset and preprocessing path.  PhysioNet 2012 source files use record identifiers, whereas MIMIC-III source tables use ICU-stay and admission identifiers.  The implemented pipelines normalize these source-specific identifiers to the canonical join key \texttt{ts\_id}.  This identifier should be read as the analysis-record identifier used by the software; it is not assumed to be a globally unique person identifier across repeated admissions or stays.

For a study unit \(i\), the observed longitudinal data are represented as a finite collection of measurement events
\[
  \mathcal{O}_i = \{(i,t_{ij},v_{ij},x_{ij}) : j=1,\ldots,n_i\},
\]
where \(i\) is the normalized \texttt{ts\_id}, \(t_{ij}\) is elapsed time in minutes from record start or ICU admission, \(v_{ij}\in\mathcal{V}\) is the observed variable name, and \(x_{ij}\) is the numeric value recorded for that variable.  In the causal repository this long-format object is stored in \texttt{ts} with the required columns \texttt{ts\_id}, \texttt{minute}, \texttt{variable}, and \texttt{value}.  Static or baseline quantities such as age, sex or gender coding, height, weight, and ICU-type indicators are represented within the same event schema as variables observed at or near the beginning of the record when the preprocessing implementation provides them.

The event representation is intentionally irregular.  Observations occur at nonuniform times, different variables are measured with different frequencies, and most variables are absent at most possible times.  A missing measurement is therefore not interpreted as a normal value.  It may reflect a clinical decision not to measure, an unrecorded observation, a source-table limitation, or a preprocessing exclusion.  The thesis treats informative measurement and missingness processes as important threats to interpretation, not as causal findings established by the repository.

The primary patient- or stay-level outcome used by the downstream causal repository is \texttt{in\_hospital\_mortality}.  In the causal processed artifact, this outcome is stored in \texttt{oc}, the outcome table keyed by \texttt{ts\_id}.  The third object in that processed artifact, \texttt{ts\_ids}, is the sorted collection of normalized identifiers represented in \texttt{ts}.  Later predictive processing uses a distinct split-aware artifact and must not be conflated with this causal contract.

\section{Design-Time Knowledge Construction and Executed Analysis}
\label{sec:problem-definition-study-design:design-time-executed-layers}

The framework distinguishes a design-time representation-construction layer from deterministic patient-level instantiation and subsequent analysis.  At design time, the official dataset-description link and instructions to research relevant clinical, causal, and time-series literature were supplied to a structured LLM-assisted elicitation protocol.  The protocol requested a dataset audit, proxy-state ontology, binary decision-rule families, reasoning about missingness and measurement processes, and a dataset-specific DAG.  The LLM received schema-level information and the mortality outcome, not patient records, empirical distributions, exposure prevalences, estimated effects, or execution access.  The resulting proposals were project-selected and encoded in deterministic tagger and graph source files.  Available evidence does not establish that every proposal was implemented unchanged, a complete line-by-line prompt-output-to-code mapping, or a formal clinician-review chain.

The analytical layer begins only after those designs have been encoded.  It preprocesses raw ICU data, applies deterministic decision-tree rules to generate rule-derived proxy labels, trains deep irregular-time-series models to predict those labels, normalizes predicted outputs, and forms the archived final aggregate from one rule-derived source and four predicted-label sources before DAG-guided matching, effect estimation, sensitivity analysis, and permutation diagnostics.  No LLM is invoked for patient-level preprocessing, label generation, prediction, adjustment selection, or effect estimation during this executed layer.  Thus prompts document design proposals, active source defines implementation, and every downstream inference remains conditional on proxy, graph, support, and model assumptions.  The two datasets follow this factorization separately; source-specific measurement availability and rule semantics are retained rather than forced into a common pooled representation.

\section{Prediction Task}
\label{sec:problem-definition-study-design:prediction-task}
% Section ID: C3.2.
% Evidence map: see thesis-writing/planning/chapter_evidence_map.md.

The prediction task is a multi-label proxy-state prediction problem over irregular ICU time series.  For each study unit \(i\) and proxy state \(k\), let \(L^{rule}_{ik}\in\{0,1\}\) denote a rule-derived proxy label.  These labels are clinically inspired analytical constructs generated from observed data by the deterministic source-code rules selected and encoded after LLM-assisted design elicitation.  They are weak clinical labels or proxy phenotypes, not chart-adjudicated diagnoses, manually verified clinical states, or validated phenotypes unless separate validation evidence is supplied in a later chapter.

The supervised predictive workflow estimates a vector of proxy-state labels from the observed time-series representation.  Its target matrix is loaded from a latent-label CSV keyed by \texttt{ts\_id}; the supervised proxy-label targets are not taken from the mortality column in \texttt{oc}.  For model \(m\), the exported prediction object may contain probabilities \(\hat p^{(m)}_{ik}\) and corresponding binary predicted labels \(\hat L^{(m)}_{ik}\).  These outputs are model estimates of rule-derived proxy labels rather than direct measurements of patient physiology.

Within construction and characterization of each testbed, five workflow functions are distinguished.  First, deterministic proxy-state construction derives binary rule-derived proxy labels from observed clinical measurements.  Second, multi-label proxy-state prediction estimates those labels from irregular time-series inputs and measures recoverability rather than clinical accuracy.  Third, the archived final causal runs aggregate one rule-derived table with four aligned binary model-prediction tables.  Fourth, mortality prediction or association analysis evaluates prognostic information in the proxy representation.  Fifth, observational causal-effect estimation treats selected aggregate columns as modeled exposures and estimates adjusted contrasts under explicit assumptions.  These functions create or characterize different resource components; none supplies known causal truth.  Detailed rules, architectures, estimators, and results follow in later chapters.

\section{Causal Question, Exposures, Outcome, Estimands, Assumptions}
\label{sec:problem-definition-study-design:causal-question-exposures-outcome-estimands-assumptions}
% Section ID: C3.3.
% Evidence map: see thesis-writing/planning/chapter_evidence_map.md.

The downstream causal-analysis setting is retrospective and observational.  No intervention is assigned by this study.  The repository uses software parameters named \texttt{treatment} for selected binary \texttt{LAT\_*} columns, but many such columns represent illness-state or organ-dysfunction proxy states rather than well-defined administered therapies.  In the thesis narrative these variables are therefore described as proxy-state exposures or modeled treatment variables unless a later section explicitly defines a manipulable intervention.

Let \(A_i\in\{0,1\}\) denote a selected proxy-state exposure, \(Y_i\) denote the binary outcome \texttt{in\_hospital\_mortality}, \(W_i\) denote observed adjustment variables selected from a project-specified DAG and available columns, and \(X_i\) denote effect-modifier features used to describe heterogeneity.  The dataset-specific DAGs were developed through a project design process that included LLM-assisted clinical and causal elicitation and were then encoded in source code; the implemented graph scripts, not the prompt outputs alone, define the DAG artifacts used by the pipeline.  Potential-outcome notation such as \(Y_i(a)\) is useful as conceptual language, but causal interpretation requires more than the presence of a graph or an estimator.  It depends on assumptions about the causal graph, consistency, exchangeability conditional on observed covariates, positivity or overlap, outcome and exposure measurement, and unmeasured confounding \cite{pearl_1995_causal_diagrams,hernan_robins_2016_target_trial,hernan_taubman_2008_well_defined_interventions}.

Accordingly, later estimates should be interpreted as adjusted effect estimates under project assumptions, not as unconditional clinical effects.  The existence of a DAG, matching script, or CATE estimator does not prove identification.  Proxy states may be useful for structured analysis while still carrying label noise, temporal ambiguity, and intervention-definition limitations.
